\scanpage{97}
\mspage{\unsure{90}}
% archive sheet number: 96
\deleted{\unsure{des $X_i$ quelconques}}
\[
X=\varinjlim_{i\in I}X_i,
\]
$I$ grand devant $\pi$, les $X_i$ dans $M$, alors $X\in M$. La condition b) \inserted{``il suffit''} l'assure$^\ast$ quand on suppose de plus les $X_i$ dans $\mathrm{Filt}^{\pi}$, (donc dans $M^{\pi}$), de sorte que l'on aura $I_M=I$, donc $I_M$ cofinal, donc $X\in M$.
% slanted
\margin{pas bien: faut \unsure{encore} supposer plus, sur le système inductif. On ne peut \unsure{tjrs} disposer \unsure{des} \unsure{lim.} \unsure{en ce cas}, $X_i\in M^{\pi}$ $\forall i\in I$, \unsure{le|et} \unsure{faire}, \unsure{le} \unsure{raisonnement} \unsure{simplement}.}
% slanted
\margin{\deleted{\uline{NB} \unsure{ce qu'il faut} \unsure{faire} \illegible{2} pour $b\Rightarrow a$.}}

Si on ne suppose pas les $X_i$ dans $\mathrm{Filt}^{\pi}$, il faut travailler un peu plus, en \deleted{prouvant ceci:}
\margin{Passer \unsure{ici} de suite à la partie de la démonstration ci-dessous page 96 (qui prouve $b\Rightarrow a$). \unsure{Le reste de la démonstr.} \unsure{n'est pas prouvé}.}
% slanted
\margin{\uuline{Ultra-faux!!}}
\deleted{\uline{!! Lemme 2} La condition b) \unsure{implique}

La condition b) \unsure{s'implique|s'entend}, plus fort, si on n'exige plus que les $X_i$ soient dans $\mathrm{Filt}^{\pi}$.

Soit en effet $x=(X_i)_{i\in I}$ avec $I$ \inserted{filtrant et} tel que tt partie filtr. $J$ de cardinal $\leq\pi$ ait une borne supérieure $i_J$, et qu'on ait $X_{i_J}\xleftarrow{\unsure{\sim|\simeq}}\varinjlim_{i\in J}X_i$. Pour chaque $i\in I$,

il existe par prop. 1', une présentation parfaite (pour $\pi$) de $X_i$,
\[
X_i=\varinjlim_{\alpha\in I_i}X_{i,\alpha}
\qquad\left\{\begin{aligned}
&\text{les }X_{i,\alpha}\in\mathrm{Filt}^{\pi}\\
&I_i\text{ grand devant }\pi\\
&\text{et de plus, tt }\unsure{A}\text{ partie de }I_i\\
&\text{cardinal }\leq\pi\text{ a une borne supérieure}\\
&\alpha_A,\text{ etc}
\end{aligned}\right.
\]}

\uline{Lemme 2} Si tout élément de $\mathrm{Filt}^{\pi}$ est $\pi$-admissible, alors $b_\pi)$ (avec $a_\pi$) implique que $M$ est stable par $\varinjlim$ grandes devant $\pi$.
